\chapter{Судалгааны арга зүй}
\label{ch:method}

%======================================================================
% Бүтэц нь яриа таних харьцуулалтын стандарт өгүүлэлтэй ижил:
% өгөгдөл → харьцуулсан систем → туршилтын явц → хэмжүүр → шалгуур.
% Жишээ: Radford нар 2023 (Whisper), Mengke нар 2025, Szymański нар 2020.
% Хайлт (RAG) ба мэдлэгийн модны арга зүйг W07–W08-ын хэмжилтийн дараа
% тусдаа хэсэг болгон нэмнэ. Хуучин текст git түүхэнд (755bf2d) бий.
%======================================================================

Энэ бүлэгт Монгол яриаг текст болгох хоёр системийг хэрхэн харьцуулсныг бичнэ. Судалгааны асуулт нь: \textit{сургуулийн тухай Монгол асуултыг текст болгоход Монгол хэлээр сургасан Whisper загвар ба Gemini 3.5 Transcribe үйлчилгээ ижил бичлэг дээр ямар алдаатай вэ?}

\section{Өгөгдөл}
\label{sec:stt-data}

Нэг хүний 13 бичлэгийг ашигласан. Бичлэг бүр сургуулийн тухай нэг асуулт. Жишээ нь элсэлт, төлбөр, сургуулийн түүхийн тухай асуулт байна. Бичлэгүүдийн нийт урт 152 секунд. Файлууд 44.1 kHz, нэг сувагтай, 16 битийн WAV.

Бичлэг бүрийн зөв бичвэрийг (reference) хүн сонсож баталсан. Зөв бичвэрт нийт 186 үг байна. Хамгийн богино бичлэг 8 үг, хамгийн урт нь 40 үгтэй.

Эдгээр бичлэгээр ямар ч загвар сургаагүй. Зөвхөн шалгахад ашигласан. Гэхдээ бичлэгүүдийг өмнө нь системийг хөгжүүлэхэд харсан. Иймээс энэ нь хөгжүүлэлтийн багц юм, тусдаа шалгалтын багц биш.

\section{Харьцуулсан системүүд}
\label{sec:stt-systems}

Хоёр системд яг ижил аудио файлыг өгсөн. Тохиргоог \tref{tab:stt_setup}-т харуулав.

\begin{table}[H]
\centering
\caption{Харьцуулсан хоёр систем}
\label{tab:stt_setup}
\begin{tabular}{|p{3.0cm}|p{5.4cm}|p{5.2cm}|}
\hline
 & \textbf{Whisper} & \textbf{Gemini 3.5 Transcribe} \\
\hline
Загвар & \texttt{Cafet/\allowbreak whisper-small-\allowbreak mongolian} \cite{cafet_whisper_small_mn} & \texttt{gemini-3.5-transcribe} \cite{google_transcribe_model} \\
\hline
Хувилбар & Hugging Face revision \texttt{d54c63a} & 2026.10.30-ны API \\
\hline
Ажилласан газар & Компьютерын CPU, faster-whisper 1.2.1, int8 & Google-ийн сервер, үнэгүй квот \\
\hline
Тохиргоо & Хэл \texttt{mn}, beam 5, temperature 0, VAD унтраасан & Хэл \texttt{mn-MN}, үгчилсэн (verbatim) горим \\
\hline
\end{tabular}
\end{table}

Whisper загвар нь OpenAI-ийн Whisper small загварыг Монгол хэлээр нэмэлт сургасан хувилбар. Загварын хуудсанд сургалтын өгөгдөл Common Voice 11 гэж бичсэн. Гэхдээ ямар хэсэг дээр, ямар дүрмээр алдааг тооцсон нь тодорхойгүй. Иймээс тэнд бичсэн алдааны хувийг энд ашиглаагүй.

Gemini-ийн загварыг бид сургах, өөрчлөх боломжгүй. Google загвараа шинэчилж болох тул бид туршилтын огноог бичсэн.

\section{Туршилтын явц}
\label{sec:stt-procedure}

Туршилтыг 2026 оны 10-р сарын 30-нд хийсэн. Бичлэг бүрийг систем бүрд нэг удаа өгсөн. Алдаа гарвал дахин илгээгээгүй. Gemini-д илгээсэн файлыг хариу авсны дараа серверээс устгасан.

13 бичлэгийн нэг нь Gemini-ийн хязгаарын алдаа (HTTP 429) авсан. Энэ нь үнэгүй квотын хязгаараас болсон. Хоёр систем хоёулаа хариу өгсөн 12 бичлэгээр харьцуулсан. Энэ 12 бичлэгт 176 үг байна.

Хугацааг бичлэг бүрээр хэмжсэн. Whisper-ийн хугацаа нь компьютер дээр текст болгосон хугацаа. Gemini-ийн хугацаанд файл илгээх, хариу авах, устгах хугацаа орсон. Иймээс хоёр хугацаа яг ижил зүйлийг хэмжээгүй.

\section{Үнэлгээний хэмжүүр}
\label{sec:stt-metrics}

Гол хэмжүүр нь үгийн алдааны хувь (\ac{WER}), томьёо (\ref{eq:wer}). Мөн тэмдэгтийн алдааны хувийг (\ac{CER}) тооцсон. CER-д хоосон зайг тоолоогүй.

Тооцохоос өмнө хоёр системийн текстэд ижил дүрэм хэрэглэсэн:
\begin{enumerate}
    \item Үсгийн хэлбэрийг Unicode NFC болгосон.
    \item Бүх үсгийг жижиг болгосон.
    \item Цэг, таслал зэрэг тэмдгийг зай болгосон.
    \item Тоог өөрчлөөгүй. Жишээ нь «2026» ба «хоёр мянга хорин зургаа» хоёрыг ялгаатай үг гэж үзсэн.
\end{enumerate}
Сүүлийн дүрэм Gemini-д хатуу. Учир нь Gemini тоог цифрээр бичих нь элбэг. Энэ асуудлыг \ref{sec:res-errors}-д дахин ярина.

Сургуулийн мэдээлэлд нэр, тоо, он чухал. Иймээс бичлэг бүрийн чухал үгсийг (нэр, тоо, он) тусад нь шалгасан.

\section{Урьдчилсан шалгуур}
\label{sec:stt-criteria}

Үр дүнг харахаас өмнө шалгуур тогтоосон. Gemini-г Whisper-ээс илүү гэж үзэхийн тулд гурван нөхцөл хангах ёстой:
\begin{enumerate}
    \item Үгийн алдаа Whisper-ээс дор хаяж 10\% бага байх.
    \item Тэмдэгтийн алдаа муудаагүй байх.
    \item Whisper зөв таньсан нэр, тоо, оныг шинээр алдаагүй байх.
\end{enumerate}

\section{Үр дүнгийн тодорхойгүй байдал}
\label{sec:stt-ci}

Бичлэг цөөн тул үр дүн санамсаргүй байдлаас их хамаарна. Үүнийг харуулахын тулд bootstrap аргаар 95\%-ийн интервал тооцсон. 12 бичлэгээс буцаан тавих замаар 12 бичлэг санамсаргүй сонгож, хоёр системийн WER-ийг тооцно. Үүнийг 10000 удаа давтана. Хоёр систем ижил бичлэгээр сонгогдоно. Тооцоог \texttt{tools/\allowbreak bootstrap\_stt\_1030.py} програмаар хийсэн.

Бүх бичлэг нэг хүнийх. Нэг хүний бичлэгүүд хоорондоо төстэй байдаг тул энэ интервал бодит тодорхойгүй байдлаас нарийн гарна. Өөрөөр хэлбэл бодит интервал үүнээс өргөн байж болно.
% ИШ ТАТАХ: Blockwise bootstrap өгүүллийг (STT15, arXiv 1912.09508) W20-д уншсаны дараа
% энд иш татна. Уншаагүй өгүүллийг иш татахгүй.
